% !TeX program = xelatex
% =====================================================================
%  Experiment 1: Object Detection and Recognition — Lab Report (English)
%  Clean single-column academic style. Compile with XeLaTeX.
% =====================================================================

\documentclass[12pt, a4paper]{article}

% ============================ Personal info (edit) ====================
\newcommand{\studentname}{Your Name}
\newcommand{\studentid}{Your Student ID}
\newcommand{\studentclass}{\begin{tabular}[t]{@{}l@{}}Class 2, Grade 2024 \\ Computer Science and Technology (Sino-Foreign Cooperative Education)\end{tabular}}
\newcommand{\teacher}{Zhou Xiaowei}
\newcommand{\expdate}{\today}

% ============================ Packages ================================
\usepackage[a4paper, top=2.5cm, bottom=2.5cm, left=2.8cm, right=2.8cm]{geometry}
\usepackage{amsmath, amssymb}
\usepackage{fontspec}
\setmainfont{TeX Gyre Termes}
\usepackage{graphicx}
\usepackage{float}
\usepackage{booktabs}
\usepackage{array}
\usepackage{setspace}
\usepackage{listings}
\usepackage{xcolor}
\usepackage{caption}
\usepackage{subcaption}
\usepackage{fancyhdr}
\usepackage{enumitem}
\usepackage{hyperref}

% ============================ Captions ================================
\captionsetup{font=small, labelfont=bf}

% ============================ Line spacing ============================
\setstretch{1.15}

% ============================ Code listings ===========================
\lstset{
    basicstyle=\ttfamily\small,
    backgroundcolor=\color{gray!12},
    numbers=left,
    numberstyle=\tiny\color{gray},
    frame=single,
    rulecolor=\color{gray},
    breaklines=true,
    captionpos=b,
    tabsize=4,
    showstringspaces=false,
}

% ============================ Header / footer =========================
\pagestyle{fancy}
\fancyhf{}
\fancyhead[L]{\small Object Detection and Recognition Lab Report}
\fancyhead[R]{\small\thepage}
\renewcommand{\headrulewidth}{0.4pt}

% ============================ Hyperlinks ==============================
\hypersetup{
    colorlinks=true,
    linkcolor=black,
    urlcolor=blue,
    citecolor=black,
    pdftitle={Object Detection and Recognition Lab Report},
    pdfauthor={\studentname},
}

% ============================ Image fallback ==========================
\newcommand{\img}[2][]{%
    \IfFileExists{#2}%
        {\includegraphics[#1]{#2}}%
        {\fbox{\parbox{0.6\linewidth}{\centering\small(Image missing: #2)}}}%
}

% ============================ Body ====================================
\begin{document}

% ============================ Cover ===================================
\begin{titlepage}
    \thispagestyle{empty}
    \vspace*{1.2cm}
    \begin{center}
        {\Large\bfseries Experiment Report}\\[1cm]
        {\LARGE\bfseries Experiment 1: Object Detection and Recognition}\\[0.5cm]
        {\large Desktop Object Detection with YOLOv8 + ROS2 + Jetson}\\[2.2cm]

        \renewcommand{\arraystretch}{1.6}
        \begin{tabular}{rl}
            \textbf{Name:} & \studentname \\[0.3em]
            \textbf{Student ID:} & \studentid \\[0.3em]
            \textbf{Class:} & \studentclass \\[0.3em]
            \textbf{Instructor:} & \teacher \\[0.3em]
            \textbf{Date:} & \expdate \\
        \end{tabular}
    \end{center}

    \vspace{1.5cm}
    {\large\bfseries Abstract}\\[0.6em]
    \noindent This report describes a desktop object detection experiment based on
    YOLOv8, deployed with ROS2 on a Jetson device. Two object classes (keyboard and
    laptop) were selected, and a dataset was independently collected and annotated,
    comprising 410 training images, 31 validation images, and 23 test objects. A
    YOLOv8n model was trained to a validation mAP50 of 0.920, and the final test
    accuracy was 87.0\%, satisfying the $\geq 80\%$ acceptance requirement. A ROS2
    node was implemented to publish detection results as
    \texttt{vision\_msgs/Detection2DArray} messages.

    \vspace{1em}
    \noindent\textbf{Keywords:} object detection; YOLOv8; ROS2; Jetson; TensorRT

    \vfill
\end{titlepage}

% ============================ TOC =====================================
\tableofcontents
\newpage

% ============================ 1. Objectives ===========================
\section{Experiment Objectives}
\begin{enumerate}[label=\arabic*., leftmargin=2em]
    \item Master the basic principles of object detection and representative algorithms (the YOLO family);
    \item Independently complete the collection, annotation, and model training of a desktop object dataset;
    \item Master the method of deploying and running a detection model in real time on Jetson edge devices;
    \item Learn to publish detection results through ROS2, integrating the detection module with a robotic system.
\end{enumerate}

% ============================ 2. Tasks ================================
\section{Experiment Tasks}
\begin{enumerate}[label=\arabic*., leftmargin=2em]
    \item Select at least 2 classes of desktop objects;
    \item Collect and annotate data independently, and train an object detection model;
    \item Run the recognition program on the Jetson;
    \item Display the object class, bounding box, and confidence in real time;
    \item Publish the recognition results via ROS2.
\end{enumerate}

% ============================ 3. Environment ==========================
\section{Experiment Environment}
\begin{table}[H]
    \centering
    \begin{tabular}{@{}lll@{}}
        \toprule
        \textbf{Category} & \textbf{Item} & \textbf{Configuration} \\
        \midrule
        Hardware & Deployment platform & NVIDIA Jetson Orin NX \\
                 & Camera & USB camera \\
        Training & CPU/GPU & Windows host (CPU training) \\
        Software & Detection framework & Ultralytics YOLOv8 \\
                 & Deep learning framework & PyTorch \\
                 & Robot middleware & ROS2 (Humble) \\
                 & Deployment acceleration & TensorRT / ONNX \\
                 & Language & Python 3 \\
        \bottomrule
    \end{tabular}
    \caption{Software and hardware environment}
\end{table}

% ============================ 4. Principles ===========================
\section{Experiment Principles}
\subsection{YOLOv8 Object Detection}
YOLOv8 adopts a one-stage detection paradigm, unifying object localization and
classification into a single regression problem. Its network consists of three parts:
Backbone, Neck, and Head, which jointly predict bounding boxes and class confidences
in a single pass over the input image. Compared with two-stage algorithms, YOLOv8
has a significant speed advantage, making it well suited for deployment on
computationally constrained edge devices such as the Jetson.

\subsection{Evaluation Metrics}
\begin{itemize}[leftmargin=2em]
    \item \textbf{IoU (Intersection over Union)}: the ratio of the intersection to the union of the predicted box and the ground-truth box, measuring localization quality;
    \item \textbf{mAP (mean Average Precision)}: the average AP over all classes, the core metric of object detection;
    \item \textbf{Accuracy}: the ratio of correctly detected objects to the total number of test objects; this experiment requires $\geq 80\%$.
\end{itemize}

% ============================ 5. Procedures ===========================
\section{Experiment Procedures}

\subsection{Dataset Collection and Annotation}
Two classes of desktop objects were selected: \texttt{keyboard} and \texttt{laptop}.
The data come from two sources: the training and validation sets use public internet
images, while the test set is entirely self-shot with a camera, ensuring that the test
set does not overlap with the training set (avoiding data leakage). After collection,
the images were annotated with LabelMe and split into training, validation, and test sets:

\begin{table}[H]
    \centering
    \begin{tabular}{@{}lccc@{}}
        \toprule
        \textbf{Dataset} & \textbf{Count} & \textbf{Source} & \textbf{Purpose} \\
        \midrule
        train & 410 images & 298 internet + 112 self-shot & Model training \\
        val   & 31 images  & internet & Validation \\
        test  & 23 objects (17 images) & all self-shot & Final evaluation \\
        \bottomrule
    \end{tabular}
    \caption{Dataset split}
\end{table}

The annotation follows the YOLO format (one object per line: class index + normalized
center coordinates and width/height).

\begin{figure}[H]
    \centering
    \begin{subfigure}{0.48\linewidth}
        \img[width=\linewidth]{labels.jpg}
        \caption{Annotation count distribution}
    \end{subfigure}
    \hfill
    \begin{subfigure}{0.48\linewidth}
        \img[width=\linewidth]{train_batch0.jpg}
        \caption{Training sample examples}
    \end{subfigure}
    \caption{Dataset annotation and samples}
\end{figure}

\subsection{Model Training}
\begin{lstlisting}[language=bash, caption={Training command}]
python train.py --epochs 60 --imgsz 480 --batch 32
\end{lstlisting}

The training parameters and results are recorded in the following table:

\begin{table}[H]
    \centering
    \begin{tabular}{@{}ll@{}}
        \toprule
        \textbf{Parameter} & \textbf{Value} \\
        \midrule
        Pretrained weights & yolov8n.pt \\
        Epochs & 60 \\
        Image size (imgsz) & 480 \\
        Batch size & 32 \\
        Initial learning rate (lr0) & 0.01 \\
        Optimizer & AdamW (auto) \\
        \bottomrule
    \end{tabular}
    \caption{Training parameters}
\end{table}

\subsection{Jetson Deployment and Acceleration}
On the Jetson, the PyTorch weights are converted to a TensorRT engine to achieve
real-time inference speed:

\begin{lstlisting}[language=bash, caption={Export TensorRT}]
yolo export model=best.pt format=engine imgsz=640
\end{lstlisting}

\subsection{ROS2 Publishing}
A ROS2 node is implemented to publish detection results as standard
\texttt{vision\_msgs/Detection2DArray} messages:

\begin{lstlisting}[language=bash, caption={Run the ROS2 node}]
source /opt/ros/humble/setup.bash
python3 ros2_detect.py
ros2 topic echo /detections
\end{lstlisting}

% ============================ 6. Results ==============================
\section{Experimental Results}

\subsection{Training Results}
\begin{figure}[H]
    \centering
    \img[width=0.95\linewidth]{results.png}
    \caption{Training loss and mAP curves}
\end{figure}

Final validation metrics: \texttt{mAP50 = 0.920, mAP50-95 = 0.734}
(epoch 60, precision 0.946 / recall 0.828).

The training curves show that the loss decreases rapidly in the early epochs and then
converges, while mAP50 rises and stabilizes, indicating that the model has converged
without obvious overfitting. Note that the precision (0.946) is higher than the recall
(0.828), meaning the model rarely produces false positives but occasionally misses
some objects, which is consistent with the error cases analyzed below.

\begin{figure}[H]
    \centering
    \img[width=0.8\linewidth]{confusion_matrix_normalized.png}
    \caption{Confusion matrix (normalized)}
\end{figure}

As shown in the normalized confusion matrix, the two classes are well separated along
the diagonal, and the main confusion is between the objects and the background, which
corresponds to the missed detections.

\subsection{Test Results and Accuracy}
\begin{table}[H]
    \centering
    \begin{tabular}{@{}p{3.5cm}cccc@{}}
        \toprule
        \textbf{Class} & \textbf{Correct} & \textbf{Total} & \textbf{Missed} & \textbf{Accuracy} \\
        \midrule
        keyboard & 6  & 8  & 2 & 75.0\% \\
        laptop   & 14 & 15 & 1 & 93.3\% \\
        \midrule
        Total    & 20 & 23 & 3 & 87.0\% \\
        \bottomrule
    \end{tabular}
    \caption{Test set results (accuracy 87.0\%, meeting $\geq 80\%$)}
\end{table}

The test set consists of 17 self-shot photos with 23 objects in total (8 keyboards
and 15 laptops), of which 20 were correctly detected, giving an overall accuracy of
87.0\%. From the per-class results, the laptop class (93.3\%) performs noticeably
better than the keyboard class (75.0\%): laptops are larger and have a more
distinctive shape, whereas keyboards often appear edge-on in the test images and
shrink into thin strips with few distinguishable features, making them harder to
detect. The figure below shows some correctly detected test samples (green boxes =
keyboard, orange boxes = laptop):

\begin{figure}[H]
    \centering
    \begin{subfigure}{0.32\linewidth}
        \img[width=\linewidth]{test1.jpg}
        \caption{Keyboard + laptop}
    \end{subfigure}
    \hfill
    \begin{subfigure}{0.32\linewidth}
        \img[width=\linewidth]{test2.jpg}
        \caption{Keyboard + laptop}
    \end{subfigure}
    \hfill
    \begin{subfigure}{0.32\linewidth}
        \img[width=\linewidth]{test3.jpg}
        \caption{Keyboard + laptop}
    \end{subfigure}

    \vspace{0.5em}

    \begin{subfigure}{0.32\linewidth}
        \img[width=\linewidth]{test4.jpg}
        \caption{Laptop}
    \end{subfigure}
    \hfill
    \begin{subfigure}{0.32\linewidth}
        \img[width=\linewidth]{test5.jpg}
        \caption{Laptop}
    \end{subfigure}
    \hfill
    \begin{subfigure}{0.32\linewidth}
        \img[width=\linewidth]{test6.jpg}
        \caption{Laptop}
    \end{subfigure}
    \caption{Test set detection results (all correctly detected)}
\end{figure}

\subsection{Real-time Detection Results}
Running \texttt{detect.py} on a PC (or the Jetson) overlays the class label, bounding
box, confidence, and FPS on the live camera feed in real time. Each detected object is
drawn with a colored bounding box (green for keyboard, orange for laptop) together with
its class name and confidence score, and the current FPS is shown in the top-left
corner. On the Jetson, the PyTorch weights are converted to a TensorRT engine to
achieve real-time performance. Some real-time detection frames captured during a live
run are shown below:

\begin{figure}[H]
    \centering
    \begin{subfigure}{0.48\linewidth}
        \img[width=\linewidth]{realtime1.png}
        \caption{Real-time detection 1}
    \end{subfigure}
    \hfill
    \begin{subfigure}{0.48\linewidth}
        \img[width=\linewidth]{realtime2.png}
        \caption{Real-time detection 2}
    \end{subfigure}

    \vspace{0.5em}

    \begin{subfigure}{0.48\linewidth}
        \img[width=\linewidth]{realtime3.png}
        \caption{Real-time detection 3}
    \end{subfigure}
    \hfill
    \begin{subfigure}{0.48\linewidth}
        \img[width=\linewidth]{realtime4.png}
        \caption{Real-time detection 4}
    \end{subfigure}
    \caption{Real-time detection frames (with class, box, confidence, and FPS)}
\end{figure}

\subsection{Error Case Analysis}
The test contains 23 objects: 20 correct, 3 missed, and 3 false positives,
producing 4 error images (below). The misses mainly arise from extreme viewing
angles (a keyboard seen edge-on, showing only a thin strip), a laptop close-up
exceeding the training distribution, and objects too near or too far such that
their scale differs greatly from the training samples. The false positives are
mostly desk shadows/textures being misjudged, or the same object being boxed
twice. The root cause is that the training set lacks samples from these extreme
viewpoints and scales, which can be improved by collecting more such data.

\begin{figure}[H]
    \centering
    \begin{subfigure}{0.48\linewidth}
        \img[width=\linewidth]{error1.jpg}
        \caption{Keyboard missed + false positive}
    \end{subfigure}
    \hfill
    \begin{subfigure}{0.48\linewidth}
        \img[width=\linewidth]{error2.jpg}
        \caption{Keyboard missed + false positive}
    \end{subfigure}

    \vspace{0.5em}

    \begin{subfigure}{0.48\linewidth}
        \img[width=\linewidth]{error3.jpg}
        \caption{Laptop missed}
    \end{subfigure}
    \hfill
    \begin{subfigure}{0.48\linewidth}
        \img[width=\linewidth]{error4.jpg}
        \caption{Laptop detected + false positive}
    \end{subfigure}
    \caption{Error cases (green = ground truth, red = prediction)}
\end{figure}

% ============================ 7. Acceptance ===========================
\section{Acceptance Criteria}
\begin{table}[H]
    \centering
    \begin{tabular}{@{}p{6.2cm}p{3.2cm}p{3.2cm}@{}}
        \toprule
        \textbf{Requirement} & \textbf{Measured Result} & \textbf{Passed} \\
        \midrule
        Recognize $\geq 2$ classes simultaneously & 2 classes (keyboard/laptop) & Yes \\
        Accuracy $\geq 80\%$ on 20 test objects & 87.0\% (20/23) & Yes \\
        Jetson real-time detection $\geq 5$ FPS & (fill in after Jetson run) & TBD \\
        Save test results and error cases & runs/test/ & Yes \\
        Publish recognition results via ROS2 & /detections & Yes \\
        \bottomrule
    \end{tabular}
    \caption{Acceptance requirements and results}
\end{table}

% ============================ 8. Code =================================
\section{Code Description}
\begin{enumerate}[label=\arabic*., leftmargin=2em]
    \item \texttt{train.py}: training script;
    \item \texttt{detector.py}: shared detector that auto-selects TensorRT/ONNX/PyTorch backend;
    \item \texttt{detect.py}: real-time detection (shows class, box, confidence, FPS);
    \item \texttt{test.py}: test evaluation (accuracy + error cases);
    \item \texttt{ros2\_detect.py}: ROS2 node (publishes Detection2DArray).
\end{enumerate}

Key code example (inference core):
\begin{lstlisting}[language=Python, caption={Detection inference}]
from ultralytics import YOLO

model = YOLO("best.engine")   # TensorRT engine
results = model(frame, conf=0.25, verbose=False)
for box in results[0].boxes:
    cls = int(box.cls[0])
    conf = float(box.conf[0])
    xyxy = box.xyxy[0].tolist()
    print(model.names[cls], conf, xyxy)
\end{lstlisting}

% ============================ 9. Run instructions =====================
\section{Run Instructions}
\begin{lstlisting}[language=bash, caption={Full workflow}]
# 1. Install dependencies
pip install -r requirements.txt

# 2. Train
python train.py

# 3. Export TensorRT on the Jetson
yolo export model=best.pt format=engine imgsz=640

# 4. Real-time detection
python detect.py

# 5. Publish via ROS2
source /opt/ros/humble/setup.bash
python3 ros2_detect.py

# 6. Test evaluation
python test.py
\end{lstlisting}

% ============================ 10. Conclusion ==========================
\section{Conclusion}
This experiment completed a desktop object detection task: two classes
(keyboard and laptop) were selected, a dataset was collected and annotated
independently (410 training images, 31 validation images, 23 test objects),
and a detection model was trained based on YOLOv8n, reaching a validation
mAP50 of 0.920. The test accuracy was 87.0\% (20/23), satisfying the
$\geq 80\%$ acceptance requirement. A ROS2 node was also implemented to
publish detection results as vision\_msgs/Detection2DArray, including
adaptation to the BoundingBox2D.center type differences across
vision\_msgs versions. The main shortcoming is that a few extreme viewpoints
(a keyboard seen edge-on, a laptop close-up) caused misses, reflecting
insufficient viewpoint and scale coverage in the training data; this can be
further improved by adding more such samples.

% ============================ Appendix ================================
\section{Appendix: GitHub Commit History}
\begin{itemize}[leftmargin=2em]
    \item Repository: \url{https://github.com/your-username/robot-object-detection}
    \item The commit history screenshot is shown below (step-by-step commits are required, not a single bulk commit):
\end{itemize}

\begin{figure}[H]
    \centering
    \img[width=0.9\linewidth]{commit_history.png}
    \caption{Git commit history screenshot (replace with your own)}
\end{figure}

\end{document}
